\documentclass{situationist-preprint}

\title{Where Moral and Prudential Reasons Diverge}
\msid{MTE-04-b}
\msfamily{MTE-04}
\msdiscipline{Metaethics}
\mstype{Consequence}
\mskeywords{why be moral; prudence; divergence; self-sacrifice; Gyges}

\begin{document}
\maketitle

\begin{abstract}
MTE-04-a answers \enquote{why be moral?} on two readings: morally, by morality itself, and prudentially, by the convergence of moral and prudential reasons over most of a life. We derive where the prudential answer fails. Prudence and morality converge because the goods that serve an agent's own life are largely goods of relationship, reputation and character, which morality protects. Divergence therefore occurs exactly where these goods are not at stake or are outweighed: in acts that are hidden and unrepeated, in costs that are large and concentrated, and in obligations whose beneficiaries cannot reciprocate. In each class the moral answer stands, and the prudential one is honestly negative.
\end{abstract}

\section{The source of convergence}\label{s:src}

By Proposition~3.4 of MTE-04-a, a situation in which the moral and the prudential ought differ is divergent, and there the prudential reading of the question receives a negative answer; by Proposition~3.5, the prudential answer fails only in such situations. Where they lie follows from why most situations converge (Assumption~2.3 of MTE-04-a). The prudential reasons that favor moral action are of three kinds: the trust and cooperation of others, which depend on reputation \citep[ch.~15]{Hobbes1651}; the relationships of friendship and love, which are partly constituted by mutual regard and fair dealing; and the agent's own character, which is shaped by what she repeatedly does \citep[II.1]{AristotleNE}. Call these the \emph{convergence goods}.

\begin{proposition}\label{p:where}
A situation is divergent only if, in it, the convergence goods are either not engaged by the choice or outweighed by a prudential cost of the moral act.
\end{proposition}

\begin{proof}
If the convergence goods are engaged and not outweighed, the prudential reasons for the moral act include them and are stronger than the prudential reasons against it, so the agent prudentially ought to do what she morally ought to do, and the situation is convergent.
\end{proof}

\section{Three classes of divergence}\label{s:classes}

Proposition~\ref{p:where} yields three classes.

\emph{Hidden and unrepeated acts.} When a wrong is undetectable and will not be repeated, neither reputation nor character is engaged. This is the case of Gyges' ring \citep[359c--360d]{PlatoRepublic}. Hume's sensible knave, who keeps to the rules except where an exception is safe, belongs here as well, though Hume doubts that such exceptions are as safe as the knave supposes \citep[IX.ii]{HumeEnquiryMorals}.

\emph{Large and concentrated costs.} When the moral act costs the agent her life, health or livelihood, as for the rescuer or the whistleblower, the convergence goods are engaged but outweighed.

\emph{Obligations without reciprocity.} Promises to the dying, duties to future generations and duties to those who can neither reward nor punish engage no reputation and no relationship that will continue.

In each class the moral reading of the question is answered exactly as in convergent cases (Theorem~3.1 of MTE-04-a): the moral reasons that make the act required are reasons to do it. Whether doing it is rational for the agent is answered by ETH-25: it is, if her deepest concerns lie beyond her own well-being. MTE-07 adds that a rational agent who judges that she morally ought to act is moved by the judgment.

\begin{objection}
Character is engaged by every choice, so no situation is divergent.
\end{objection}

\begin{reply}
A single hidden act shapes character little, and the cost in the second class can exceed whatever it does. To deny divergence is to deny the plain verdict about the rescuer who dies, which MTE-04-a rightly accepts.
\end{reply}

\section{Conclusion}

Divergence is located where the goods that make morality prudent are absent or outweighed. That these classes are narrow is what bounds the failures of the prudential answer; that they are not empty is why the moral answer is needed.

\begin{thebibliography}{99}\small
\bibitem[Aristotle()]{AristotleNE} Aristotle. \emph{Nicomachean Ethics}. Cited by book and chapter.
\bibitem[Hobbes(1651)]{Hobbes1651} Hobbes, T. (1651). \emph{Leviathan}. Andrew Crooke.
\bibitem[Hume(1751)]{HumeEnquiryMorals} Hume, D. (1751). \emph{An Enquiry Concerning the Principles of Morals}. Cited by section and part.
\bibitem[Plato()]{PlatoRepublic} Plato. \emph{Republic}. Cited by Stephanus pagination.
\end{thebibliography}

\end{document}
